\begin{solapartat}
\punts{0,2}
\[ \left.\begin{aligned}
F_g &= ma_c\\
G\,\frac{M_J m_\mathrm{sonda}}{r^2} &= m_\mathrm{sonda}\,\frac{v^2}{r}\\
E_c &= \frac{1}{2}m_\mathrm{sonda}v^2\\
r &= R_{\mathit{Jupiter}} + h
\end{aligned}\right\}\Rightarrow E_c = \frac{1}{2}G\,\frac{M_J m_\mathrm{sonda}}{r} \]

\punts{0,1} $r = 69911 + 4300 = 74\,211\ \mathrm{km} = 7,42\times 10^{7}\ \mathrm{m}$

\punts{0,3} $E_c = \dfrac{1}{2}G\,\dfrac{M_J m_\mathrm{sonda}}{r} = 3,10\times 10^{12}\ \mathrm{J}$

\textbf{El valor de la velocitat de Juno en la seva òrbita és:}

\punts{0,2} $v = \sqrt{\dfrac{2E_c}{m_\mathrm{sonda}}} = 4,13\times 10^{4}\ \mathrm{m\,s^{-1}}$

\punts{0,2} $T = \dfrac{2\pi r}{v} = \dfrac{2\pi\cdot 7,4211\times 10^{7}}{4,132\times 10^{4}} = 11300\ \mathrm{s}$
\end{solapartat}

\begin{solapartat}
\punts{0,4} $E_\text{òrbita} + E_\text{a comunicar} = 0$

\punts{0,6} $E_\text{a comunicar} = G\,\dfrac{M_J m_\mathrm{sonda}}{2r} = 3,10\times 10^{12}\ \mathrm{J}$
\end{solapartat}
